\chapter{Bäume und ihre Empfindlichkeit}

\section{Der Aufbau}

\begin{merksatz}
\boxtitel{Zwei Arten von Knoten}
\begin{center}
{\small
\renewcommand{\arraystretch}{1.55}
\begin{tabular}{@{}P{3.4cm}P{5.2cm}@{}}
\toprule
\textbf{Entscheidungs\-knoten} & Man wählt --- gezeichnet als Quadrat. Auswertung: das Beste nehmen.\\
\textbf{Zufallsknoten} & Die Natur wählt --- gezeichnet als Kreis. Auswertung: mit Wahrscheinlichkeiten mitteln.\\
\bottomrule
\end{tabular}}
\end{center}

\schritt{Die Auswertung}
Man beginnt an den Blättern und rechnet rückwärts zur Wurzel --- \Begriff{Rückwärtsinduktion}. Dasselbe Verfahren trägt Kapitel~7 und, in anderer Gestalt, den Spielbaum aus Band 21, Kapitel~6.

\begin{center}
\begin{minipage}{0.88\textwidth}
\kursiv{\textbf{Der Unterschied zum Spielbaum: Dort maximiert der Gegenspieler seinen eigenen Nutzen, hier würfelt die Natur.} An die Stelle des Minimums tritt der Erwartungswert. Wer die Natur als Gegner behandelt und das Schlechteste annimmt, rechnet nach Kapitel~11 --- und kommt regelmäßig zu anderen Ergebnissen.}
\end{minipage}
\end{center}
\end{merksatz}

\begin{merksatz}
\boxtitel{Der Leitfall als Baum}
\begin{center}
\begin{tikzpicture}[scale=1.0, every node/.style={font=\scriptsize}]
\node[draw=tief, fill=tief!10, rounded corners=1pt, inner sep=3pt] (w) at (0,0) {};
\node[knotenr] (a) at (2.2,1.2) {};
\node[knotenr] (b) at (2.2,-1.2) {};
\draw[kanted, -{Stealth[length=4pt]}] (w)--(a) node[midway,above,color=grau] {behandeln};
\draw[kanted, -{Stealth[length=4pt]}] (w)--(b) node[midway,below,color=grau] {nicht};
\foreach \y/\t/\v in {1.9/{Schlaganfall $0{,}02$}/{$0{,}40$}, 1.2/{Blutung $0{,}03$}/{$0{,}70$}, 0.5/{nichts $0{,}95$}/{$1{,}00$}}
  {\draw[kanted] (a)--(4.0,\y); \node[anchor=west, color=grau] at (4.05,\y) {\t\ $\to$ \v};}
\foreach \y/\t/\v in {-0.5/{Schlaganfall $0{,}06$}/{$0{,}40$}, -1.2/{Blutung $0{,}01$}/{$0{,}70$}, -1.9/{nichts $0{,}93$}/{$1{,}00$}}
  {\draw[kanted] (b)--(4.0,\y); \node[anchor=west, color=grau] at (4.05,\y) {\t\ $\to$ \v};}
\node[color=tief] at (2.2,1.75) {$0{,}9790$};
\node[color=grau] at (2.2,-1.75) {$0{,}9610$};
\end{tikzpicture}
\end{center}

\begin{center}
\begin{minipage}{0.9\textwidth}
\textbf{Der Baum macht sichtbar, was die Rechnung aus Kapitel~1 verbarg: jede Zahl steht an genau einer Kante.}
\end{minipage}
\end{center}
\end{merksatz}

\section{Empfindlichkeit}

\begin{merksatz}
\boxtitel{Drei Prüfungen}
\begin{center}
{\small
\renewcommand{\arraystretch}{1.55}
\begin{tabular}{@{}P{2.6cm}P{6.0cm}@{}}
\toprule
\textbf{Einfach} & Eine Größe über ihren plausiblen Bereich verändern; die übrigen festhalten. Ergebnis: ein Kipppunkt oder keiner.\\
\textbf{Wirbelsturm} & Alle Größen einzeln durchspielen und die Wirkungen der Größe nach ordnen --- die längsten Balken oben.\\
\textbf{Wahrscheinlich\-keitsverteilt} & Alle Größen zugleich aus Verteilungen ziehen und die Häufigkeit der Entscheidungen auszählen.\\
\bottomrule
\end{tabular}}
\end{center}

\schritt{Am Leitfall: die einfache Prüfung}
\begin{center}
{\small
\renewcommand{\arraystretch}{1.46}
\begin{tabular}{@{}P{2.8cm}P{2.2cm}P{2.6cm}@{}}
\toprule
\textbf{Größe} & \textbf{Kipppunkt} & \textbf{plausibel?}\\
\midrule
Schlaganfallrisiko & $1{,}5\,\%$ & ja, bei jüngeren\\
Blutungsrisiko mit & $9{,}0\,\%$ & selten\\
Nutzen Blutung & $-0{,}20$ & \textcolor{gruen}{nein}\\
Wirksamkeit & $r > 0{,}75$ & \textcolor{gruen}{nein}\\
\bottomrule
\end{tabular}}
\end{center}

\begin{center}
\begin{minipage}{0.9\textwidth}
\textbf{Nur eine einzige Größe kann die Entscheidung im plausiblen Bereich kippen: das Schlaganfallrisiko.}
\end{minipage}
\end{center}

\begin{center}
\begin{minipage}{0.88\textwidth}
\kursiv{\textbf{Damit ist die ganze Analyse auf eine Frage zusammengeschrumpft} --- und genau das ist ihr Zweck. Man muss nicht alle Zahlen genau kennen, sondern nur die eine, an der es hängt. Bei Frau B.\ ist sie bekannt, also ist die Entscheidung robust; bei einem $60$-jährigen Patienten ohne weitere Risikofaktoren wäre sie es nicht.}
\end{minipage}
\end{center}
\end{merksatz}

\begin{merksatz}
\boxtitel{Die Zweigrößenprüfung}
Schlaganfallrisiko $s$ und Blutungsrisiko unter Behandlung $b$ zugleich:
\[ \text{behandeln, wenn} \quad \tfrac23\,s\cdot0{,}60 \ > \ (b-0{,}01)\cdot0{,}30 \]
\[ \Longleftrightarrow \quad s \ > \ 0{,}75\,(b - 0{,}01) \]

\begin{center}
{\small
\renewcommand{\arraystretch}{1.46}
\begin{tabular}{@{}P{2.4cm}P{2.6cm}P{2.6cm}@{}}
\toprule
$b$ & \textbf{Grenze für } $s$ & \textbf{Frau B. ($6\,\%$)}\\
\midrule
$0{,}03$ & $1{,}5\,\%$ & \textcolor{gruen}{behandeln}\\
$0{,}06$ & $3{,}8\,\%$ & \textcolor{gruen}{behandeln}\\
$0{,}09$ & $6{,}0\,\%$ & Grenzfall\\
$0{,}12$ & $8{,}3\,\%$ & \textcolor{akzent}{nicht}\\
\bottomrule
\end{tabular}}
\end{center}

\begin{center}
\begin{minipage}{0.88\textwidth}
\kursiv{\textbf{Die Grenze ist eine Gerade in der Ebene der beiden Risiken}, und die Entscheidung hängt nur davon ab, auf welcher Seite man liegt. Klinische Punktwerte für Schlaganfall- und Blutungsrisiko liefern genau die beiden Koordinaten --- ihre gemeinsame Auswertung ist also nichts anderes als das Ablesen dieser Geraden.}
\end{minipage}
\end{center}
\end{merksatz}

\section{Der Bruchfall}

\begin{warnung}
\boxtitelr{Was die Prüfung nicht prüft}
\begin{center}
{\small
\renewcommand{\arraystretch}{1.52}
\begin{tabular}{@{}P{3.2cm}P{5.4cm}@{}}
\toprule
\textbf{Der Aufbau} & Ein vergessener Ast wird nie geprüft. Die Empfindlichkeit betrifft Zahlen, nicht die Struktur.\\
\textbf{Abhängigkeiten} & Größen werden meist unabhängig gezogen --- Schlaganfall- und Blutungsrisiko hängen aber zusammen.\\
\textbf{Scheingenauigkeit} & Eine Verteilungsrechnung mit erfundenen Verteilungen erzeugt Vertrauensintervalle, die nichts bedeuten.\\
\bottomrule
\end{tabular}}
\end{center}

\begin{center}
\begin{minipage}{0.88\textwidth}
\kursiv{Die zweite Zeile ist der häufigste Fehler in veröffentlichten Modellen. \textbf{Wer ein hohes Schlaganfallrisiko hat, hat meist auch ein hohes Blutungsrisiko} --- beide steigen mit Alter und Gebrechlichkeit. Zieht man sie unabhängig, so erzeugt man Patienten, die es nicht gibt, und unterschätzt die Streuung des Ergebnisses.}
\end{minipage}
\end{center}
\end{warnung}

\section{Übungsaufgaben}

\begin{aufgabe}
\boxtitel{Aufgabe 5.1 --- Einen Kipppunkt finden}
Bei welchem Blutungsrisiko unter Behandlung kippt die Empfehlung für Frau B.?

\loesung
\schritt{Schritt 1 --- Die Gleichung}
\[ \tfrac23\cdot0{,}06\cdot0{,}60 = (b-0{,}01)\cdot0{,}30 \]

\schritt{Schritt 2 --- Auflösen}
\[ 0{,}0240 = 0{,}30\,b - 0{,}0030 \quad\Longrightarrow\quad b = \frac{0{,}0270}{0{,}30} = 0{,}090 \]

\schritt{Schritt 3 --- Die Probe}
\[ \text{mit } b = 0{,}09: \quad 0{,}02\cdot0{,}40 + 0{,}09\cdot0{,}70 + 0{,}89 = 0{,}0080+0{,}0630+0{,}8900 = 0{,}9610 \]
\kursiv{Genau der Nutzen ohne Behandlung.} \checkmark

\begin{center}
\begin{minipage}{0.88\textwidth}
\kursiv{\textbf{Das Blutungsrisiko müsste auf neun Prozent steigen --- das Dreifache des angenommenen.} Solche Werte kommen vor, etwa bei wiederholten Stürzen oder schwerer Niereninsuffizienz; dann ist die Entscheidung tatsächlich offen. Der Kipppunkt sagt also, worauf man bei der Anamnese achten muss.}
\end{minipage}
\end{center}
\end{aufgabe}

\begin{aufgabe}
\boxtitel{Aufgabe 5.2 --- Der Wirbelsturm}
Ordnen Sie vier Größen nach ihrer Wirkung auf den Nutzengewinn, wenn jede um $\pm50$ Prozent schwankt.

\loesung
\schritt{Schritt 1 --- Die Rechnungen}
Gewinn $= \tfrac23 s\,\Delta u - (b-b_0)(1-u_B)$, Ausgangswert $+0{,}0180$.

\begin{center}
{\small
\renewcommand{\arraystretch}{1.46}
\begin{tabular}{@{}P{2.4cm}P{1.8cm}P{1.8cm}P{2.0cm}@{}}
\toprule
\textbf{Größe} & $-50\,\%$ & $+50\,\%$ & \textbf{Spanne}\\
\midrule
$s$ & $+0{,}0060$ & $+0{,}0300$ & $\mathbf{0{,}0240}$\\
$b$ & $+0{,}0225$ & $+0{,}0135$ & $0{,}0090$\\
$\Delta u$ & $+0{,}0060$ & $+0{,}0300$ & $\mathbf{0{,}0240}$\\
$1-u_B$ & $+0{,}0225$ & $+0{,}0135$ & $0{,}0090$\\
\bottomrule
\end{tabular}}
\end{center}

\schritt{Schritt 2 --- Die Rangfolge}
\begin{center}
\begin{minipage}{0.9\textwidth}
\textbf{Schlaganfallrisiko und Schlaganfallschwere wirken gleich stark --- und beide dreimal so stark wie die Blutungsgrößen.}
\end{minipage}
\end{center}

\schritt{Schritt 3 --- Warum das so ist}
\begin{center}
\begin{minipage}{0.88\textwidth}
\kursiv{Beide gehen als Produkt in denselben Term ein; eine Verdopplung des einen wirkt wie eine des anderen. \textbf{Das ist eine allgemeine Eigenschaft: In multiplikativen Termen ist die relative Empfindlichkeit für alle Faktoren gleich.} Wer also das Schlaganfallrisiko genauer bestimmen will, sollte ebenso viel Mühe auf die Frage verwenden, wie schwer die verhinderten Schlaganfälle wären --- was in der Praxis fast nie geschieht.}
\end{minipage}
\end{center}
\end{aufgabe}

\begin{aufgabe}
\boxtitel{Aufgabe 5.3 --- Ein vergessener Ast}
Das Modell übersieht, dass ein Teil der Schlaganfälle unter Behandlung als Hirnblutung auftritt --- also schwerer verläuft. Wie ändert sich das Ergebnis?

\loesung
\schritt{Schritt 1 --- Der neue Ast}
Von den zwei Prozent Schlaganfällen unter Behandlung sei ein Viertel eine Hirnblutung mit Nutzen $0{,}20$ statt $0{,}40$.

\schritt{Schritt 2 --- Neu rechnen}
\[ \text{mit: } 0{,}005\cdot0{,}20 + 0{,}015\cdot0{,}40 + 0{,}03\cdot0{,}70 + 0{,}95 \]
\[ = 0{,}0010 + 0{,}0060 + 0{,}0210 + 0{,}9500 = 0{,}9780 \]

\schritt{Schritt 3 --- Der Vergleich}
\begin{center}
{\small
\renewcommand{\arraystretch}{1.48}
\begin{tabular}{@{}P{2.8cm}P{2.2cm}P{2.4cm}@{}}
\toprule
& \textbf{Nutzen} & \textbf{Gewinn}\\
\midrule
altes Modell & $0{,}9790$ & $+0{,}0180$\\
mit neuem Ast & $0{,}9780$ & $+0{,}0170$\\
\bottomrule
\end{tabular}}
\end{center}

\begin{center}
\begin{minipage}{0.9\textwidth}
\textbf{Die Empfehlung ändert sich nicht --- aber keine Empfindlichkeitsprüfung des alten Modells hätte das zeigen können.}
\end{minipage}
\end{center}

\schritt{Schritt 4 --- Die Lehre}
\begin{center}
\begin{minipage}{0.88\textwidth}
\kursiv{\textbf{Strukturelle Unsicherheit ist die unbehandelte Unsicherheit der Entscheidungsanalyse.} Man begegnet ihr, indem man mehrere Modellvarianten rechnet und ihre Ergebnisse nebeneinanderstellt --- nicht, indem man Verteilungen feiner wählt. Das entspricht genau der Lage aus Band 23, Kapitel~4: Der Graph steht nicht in den Daten, und das Modell steht nicht im Baum.}
\end{minipage}
\end{center}
\end{aufgabe}

\begin{geschichte}
\boxtitelg{Vom Ölfeld ans Krankenbett}
Die Entscheidungsbäume kamen aus der Betriebswirtschaft --- \Person{Howard Raiffa} und \Person{Robert Schlaifer} lehrten sie ab den 1950er Jahren in Harvard.

\begin{center}
\begin{minipage}{0.9\textwidth}
\kursiv{\textbf{Der erste medizinische Entscheidungsbaum wurde 1967 von \Person{Lee Lusted} veröffentlicht}, der zugleich die Fachgesellschaft für medizinische Entscheidungsfindung gründete. Ihre Zeitschrift \kursiv{Medical Decision Making} erscheint seit 1981 und ist bis heute das Sammelbecken dieses Gebiets.}
\end{minipage}
\end{center}

\schritt{Und ein berühmter Baum}
\begin{center}
\begin{minipage}{0.9\textwidth}
\kursiv{1978 veröffentlichten Pauker und Mitarbeiter eine Analyse zur Frage, ob bei einem Kind mit Fieber unklarer Ursache eine Lumbalpunktion erfolgen soll. \textbf{Das Ergebnis widersprach der damaligen Praxis} und wurde heftig angegriffen --- nicht wegen der Rechnung, sondern wegen der Nutzenwerte, die die Verfasser gewählt hatten. Der Streit zeigte beispielhaft, was Kapitel~1 behauptet: Man streitet über Werte, nicht über Zahlen.}
\end{minipage}
\end{center}
\end{geschichte}

\begin{forschung}
\boxtitelp{Modelle heute}
\begin{center}
{\small
\renewcommand{\arraystretch}{1.45}
\begin{tabular}{@{}P{3.4cm}P{4.8cm}@{}}
\toprule
\textbf{Wertorientierte Prüfung} & Empfindlichkeit nicht gegen Zahlen, sondern gegen Nutzenwerte.\\
\textbf{Strukturprüfung} & Mehrere Modellformen parallel, mit Gewichtung.\\
\textbf{Offene Modelle} & Veröffentlichte Rechenwege statt Ergebnistabellen.\\
\bottomrule
\end{tabular}}
\end{center}

\schritt{Zur dritten Zeile}
\begin{center}
\begin{minipage}{0.88\textwidth}
\kursiv{Gesundheitsökonomische Modelle wurden lange als Betriebsgeheimnis behandelt --- und ihre Ergebnisse hingen, wie Nachrechnungen zeigten, oft an wenigen unveröffentlichten Annahmen. \textbf{Mehrere Zeitschriften verlangen inzwischen den vollständigen Code}, und einige Behörden rechnen Modelle unabhängig nach. Das ist dieselbe Entwicklung, die Band 22 für die Lernverfahren beschreibt: Die Prüfbarkeit wird wichtiger genommen als die Rechnung.}
\end{minipage}
\end{center}
\end{forschung}

\begin{merksatz}
\boxtitel{Was Sie jetzt können}
\begin{center}
\renewcommand{\arraystretch}{1.5}
\begin{tabular}{@{}P{3.6cm}P{5.0cm}@{}}
\toprule
\textbf{Zwei Knoten} & wählen und würfeln\\
\textbf{Rückwärtsinduktion} & von den Blättern zur Wurzel\\
\textbf{Kipppunkt} & wo die Entscheidung wechselt\\
\textbf{Wirbelsturm} & welche Größe trägt\\
\textbf{Der Bruchfall} & die Struktur wird nie geprüft\\
\bottomrule
\end{tabular}
\end{center}

\begin{center}
\begin{minipage}{0.88\textwidth}
\kursiv{Ein Baum bildet einen Zeitpunkt ab. Das nächste Kapitel verfolgt Frau B.\ über zwanzig Jahre --- und muss dafür Zustände wiederholen, statt Äste zu vervielfachen.}
\end{minipage}
\end{center}
\end{merksatz}
